\section{Архитектура разработанного микросервиса}

\subsection{Общая архитектурная идея}

Центральный архитектурный результат -- не столько выбор конкретного
дерева, сколько выделение сопоставления «префикс~$\rightarrow$~ASN» в
отдельный сетевой процесс между источником BGP и продуктами компании.
Сервис занимает промежуточное положение между узлом GoBGP и
потребителями; последние подключаются по одному и тому же gRPC контракту,
независимо от того, реализован backend на \texttt{ranger} или на
шардированном ART. В качестве источника используется GoBGP с gRPC API;
в качестве потребителей -- BIFIT MITIGATOR и BIFIT COLLECTOR. Логически
архитектура изображена на рис.~\ref{fig:arch-top}.

\begin{figure}[ht]
\centering
\begin{tikzpicture}[node distance=1.6cm and 1.8cm]
\node[block, text width=2.8cm] (cl)  {Потребители\\ \textit{Mitigator}, \textit{Collector}};
\node[block, text width=3.1cm, right=of cl] (svc) {\textbf{Сервис}\\ сопоставления};
\node[block, text width=2.9cm, right=of svc] (bgp) {\textbf{GoBGP}\\ маршрутная база};

\draw[msarr] (cl.east)  -- node[above,font=\footnotesize]{gRPC} (svc.west);
\draw[msarr] (svc.east) -- node[above,font=\footnotesize]{gRPC} (bgp.west);
\end{tikzpicture}
\caption{Общая архитектура: сервис между потребителями и узлом GoBGP}
\label{fig:arch-top}
\end{figure}

На логическом уровне внутри сервиса можно выделить следующие
компоненты:

\begin{itemize}
\item модуль конфигурации и запуска;
\item gRPC-сервер пользовательского API;
\item клиент GoBGP;
\item доменный сервис синхронизации состояния;
\item хранилище префиксов \texttt{prefixStore};
\item подсистема рассылки обновлений подписчикам (subscriber hub);
\item вспомогательные механизмы диагностики и профилирования (pprof).
\end{itemize}

Такое разделение важно, потому что оно изолирует два разных типа
ответственности: получение и актуализацию маршрутизационного состояния
с одной стороны и обслуживание внешних запросов с другой. Внутренние
логические слои показаны на рис.~\ref{fig:arch-layers}.

\begin{figure}[ht]
\centering
\begin{tikzpicture}[node distance=0.45cm,every node/.style={font=\footnotesize}]
\node[block, fill=gray!10, text width=11cm] (entry) {Точка входа и конфигурация};
\node[block, fill=gray!10, text width=11cm, below=of entry] (svc)
   {Доменный сервис синхронизации (lifecycle: connect $\to$ snapshot $\to$ load $\to$ watch $\to$ resync)};
\node[block, fill=gray!10, text width=5.3cm, below=of svc, xshift=-3cm] (bgp)
   {Клиент GoBGP\\ \texttt{ListPath}, \texttt{WatchEvent}};
\node[block, fill=gray!10, text width=5.3cm, right=of bgp] (store)
   {Хранилище \texttt{prefixStore}\\ (ranger / шард. ART)};
\node[block, fill=gray!10, text width=11cm, below=0.6cm of bgp.south, anchor=north, xshift=3cm] (api)
   {gRPC-сервер пользовательского API + subscriber hub};

\draw[msarr] (entry.south) -- (svc.north);
\draw[msarr] (svc.south) -- (bgp.north);
\draw[msarr] (svc.south) -- (store.north);
\draw[msarr] (store.south) -- (api.north);
\draw[msarr] (bgp.south) -- (api.north);
\end{tikzpicture}
\caption{Внутренние логические слои сервиса}
\label{fig:arch-layers}
\end{figure}

\subsection{Модель взаимодействия с GoBGP}

С GoBGP сервис говорит по двум каналам.

Первый канал, \texttt{ListPath}, отдаёт полный снимок таблицы. Это
операция начальной загрузки, но при необходимости она вызывается
повторно -- например, при пересинхронизации.

Второй канал, \texttt{WatchEvent}, поставляет поток событий: новые
анонсы и отзывы маршрутов в том порядке, в котором их видит соседний
BGP-роутер.

Получается жизненный цикл такого вида:

\begin{enumerate}
\item подключиться к GoBGP;
\item прочитать полный снимок;
\item нормализовать и загрузить его во внутреннее хранилище;
\item перейти в режим обработки потока обновлений;
\item при ошибке или разрыве -- переподключиться и пересинхронизироваться.
\end{enumerate}

По смыслу это та же схема «начальная загрузка от снимка плюс
непрерывный поток», что и в распределённых системах обработки
событий~\cite{akidau2015dataflow,carbone2015flink}.

\subsection{Пользовательский gRPC API}

Для внешних потребителей реализуются три основных операции.

\texttt{GetASNInfo} выполняет lookup по IP-адресу и возвращает
информацию о найденном ASN и наиболее специфичном префиксе.

\texttt{GetPrefixesASTable} возвращает полный снимок таблицы
соответствия префиксов автономным системам.

\texttt{SubscribePrefixASUpdates} позволяет клиенту получать
непрерывный поток изменений после момента подписки.

Такой набор операций покрывает как synchronous query-сценарии, так и
reactive-сценарии непрерывного наблюдения за изменениями. Подробное
обсуждение семантики этих методов и поведения при отказах вынесено в
приложение~В.

\subsection{Два режима хранения: legacy и новый}

В процессе работы сервиса поддерживаются два варианта backend-хранилища
(см. рис.~\ref{fig:store-contract}):

\begin{itemize}
\item \textbf{ranger} -- прежнее решение, оптимизированное под полную
замену таблицы;
\item \textbf{шардированный ART} -- новое решение, поддерживающее
инкрементальные модификации.
\end{itemize}

Наличие двух вариантов полезно не только для постепенной миграции, но
и для экспериментов: оно позволяет на одном и том же сервисном
контуре измерять компромиссы между batch-оптимизированным и
stream-ориентированным подходами в идентичных условиях.

\begin{figure}[ht]
\centering
\begin{tikzpicture}[node distance=0.55cm,every node/.style={font=\footnotesize}]
\node[block, text width=10cm, fill=gray!12] (iface)
   {\textbf{Абстракция \texttt{prefixStore}}\\
    \texttt{Replace} $|$ \texttt{Upsert} $|$ \texttt{Delete} $|$
    \texttt{Lookup} $|$ \texttt{Dump}};
\node[block, text width=4.6cm, below left=1cm and 0.4cm of iface.south, fill=gray!4] (rg)
   {\textbf{Реализация \textit{ranger}}\\
    плоский слайс\\ полная замена при \texttt{Upsert}};
\node[block, text width=4.6cm, below right=1cm and 0.4cm of iface.south, fill=gray!4] (art)
   {\textbf{Реализация \textit{шард. ART}}\\
    шарды + ART-деревья\\ локальные \texttt{Upsert}/\texttt{Delete}};

\draw[msarr] (iface.south) -- (rg.north);
\draw[msarr] (iface.south) -- (art.north);
\end{tikzpicture}
\caption{Контракт \texttt{prefixStore} и две реализации}
\label{fig:store-contract}
\end{figure}

\subsection{Эксплуатационные ограничения}

При проектировании микросервиса с самого начала фиксируются
несколько эксплуатационных ограничений, которые дальше определяют
устройство внутренних механизмов.

\begin{itemize}
\item \textbf{Один источник маршрутных данных.} Сервис подключается к
одному узлу GoBGP. Это сознательное упрощение: оно делает модель
состояния однозначной и снимает вопрос о конфликтном слиянии данных
от нескольких BGP-роутеров. Архитектура расширению не препятствует,
но текущая версия работает по правилу «один источник -- один поток
событий -- одно внутреннее состояние».
\item \textbf{Одна копия состояния на всех клиентов.} Сервис держит
одну логическую копию маршрутной таблицы. Альтернативная схема
с копией на каждого клиента потребовала бы либо лишнего расхода
памяти, либо более сложной сериализации обновлений; ни то ни другое
для рассматриваемых сценариев не оправдано.
\item \textbf{Гарантии видимости.} После того как обновление
применилось к хранилищу, оно становится видимым последующим
читателям и подписчикам. Строгой глобальной упорядоченности видимости
для разных читателей не требуется -- это согласуется с моделью BGP,
где распределённое состояние сходится во времени.
\item \textbf{Время отклика.} \texttt{Lookup} остаётся синхронным
вызовом в пределах одного gRPC-запроса. Тяжёлые операции
(\texttt{GetPrefixesASTable}, подписка на поток) живут в отдельных
RPC и не блокируют точечные запросы.
\end{itemize}

\subsection{Вынесение логики в отдельный микросервис как основной архитектурный результат}

До внедрения отдельного сервиса таблица «префикс~$\rightarrow$~ASN»
строилась и обновлялась внутри BIFIT COLLECTOR: второй продукт не мог
опереться на ту же реализацию без копирования кода и состояния, а
смена backend потребовала бы правок в прикладной логике коллектора.
После выделения функции в микросервис MITIGATOR и COLLECTOR становятся
клиентами одного контракта: граница процесса совпадает с границей
предметной задачи -- сопоставление адресов и префиксов ASN по данным
BGP.

Такое разделение даёт и эксплуатационный эффект. Обновление алгоритма
загрузки снимка, политики нормализации префиксов или типа хранилища
(\texttt{ranger} против шардированного ART) локализуется в одном месте;
потребители продолжают вызывать те же RPC. Согласованность данных между
продуктами перестаёт зависеть от дублирования кода; сопровождение не
распадается на две независимые реализации одной и той же таблицы.

%==============================================================
